\documentclass[11pt]{article}
\usepackage{amsmath,amssymb,geometry}
\geometry{margin=1in}
\title{A counterexample to Conjecture 00000008541}
\author{gaochengzhecpu}
\date{}
\begin{document}
\maketitle

The free lattice on two generators is finite: it has precisely the four
elements
\[m=x\wedge y,\quad x,\quad y,\quad j=x\vee y.\]
Here $m<x<j$ and $m<y<j$, with $x$ and $y$ incomparable. Consequently every
set of elements of length at most $\ell$ has at most four elements, for
every $\ell$, whatever ordinary choice of length is made. Its growth is
bounded and therefore cannot be doubly exponential.

\paragraph{Verification of freeness.}
Let $A$ be an arbitrary lattice and let $a,b\in A$. There is a homomorphism
from the four-element lattice to $A$ given by
\[
 m\mapsto a\wedge b,\qquad x\mapsto a,\qquad
 y\mapsto b,\qquad j\mapsto a\vee b.
\]
For comparable pairs the meet and join identities follow immediately from
their order; for the incomparable pair $x,y$ they follow from the displayed
assignments. Thus this map preserves both operations. It is the unique
homomorphism sending $x$ to $a$ and $y$ to $b$, because the images of
$m=x\wedge y$ and $j=x\vee y$ are forced. This is the full universal property
of the free lattice on two generators, in the category of all lattices, not
merely distributive lattices. It also proves that no additional elements
or identifications were introduced in choosing the four-element model.

The source imposes no lower restriction such as $n\ge3$. This counterexample
addresses its written statement, at $n=2$. It makes no assertion about a
version restricted to larger numbers of generators. The lattice signature
here is the usual pair of binary operations $\wedge,\vee$, with no additional
nullary constants, as in the source's unconstrained lattice identities.

\paragraph{Lean verification.}
The structure \texttt{Lattice} contains an actual partial order and the
greatest-lower-bound and least-upper-bound axioms. The finite four-element
model satisfies all of them. The theorem
\texttt{free\_on\_two\_generators} proves existence and uniqueness of the
homomorphism above into every such target lattice.

The theorem \texttt{conjecture8541\_counterexample} quantifies over every
family of predicates on all four elements, so it includes any possible
family of length balls. It proves that every count is at most four and
that the count function is not unbounded. Every doubly exponential growth
law entails unboundedness; thus no convention for the unspecified word
length can rescue the claimed growth at $n=2$.

All finite decisions are checked by Lean's kernel, using Lean 4.19.0 and
\texttt{Std}. No admitted lemmas, native evaluation proofs, or custom axioms
are used.
\end{document}
